\subsection{The nonunit-inner ordinary eighth moment}\label{proof:Q5-C075}
\paragraph{Q75.0 Exact new theorem and its original scope}

Let the original finite probability space and original real skew kernel satisfy
the conventions in \mbox{\hyperref[part:Q75.1]{part Q75.1}}. Put
\fitdisplay{
W_s=A \bar f_{s-1},\qquad X_{b,L}=T^L(qf_b).
}{}
The claim is
\fitdisplay{
\boxed{E(W_sX_{b,L})^8\le e^{-3d/20}\le1-d/7,
 \qquad s\ge3,\quad b\ge2,\quad L\ge2.}
}{Q75.T080.1}
In fact the shortest new arm has the stronger bound
\fitdisplay{
E(W_3X_{b,L})^8\le e^{-d}\qquad(b\ge2,L\ge2).
}{Q75.T080.2}


At an actual root with distinct residual arms (s, a), \(1\le s<a,\allowbreak \) the negative
canonical factor is \(W_s \bar f_a.\) The short arm's first edge points away from
the root, and all its remaining edges point toward the root. The long arm
points entirely toward the root. At the terminal with residual arms (1, b),
the unit arm points toward the terminal and its b-arm points away. The
inner connector of length L points from terminal to root. These descriptions
are literal directions of every original subdivision edge.

After multiplication by the genuine equal-pair decorations \(E_r,\allowbreak  E_t,\allowbreak \) the
actual input and actual receiving-root output are
\fitdisplay{
Y_D=E_r W_sT^L(E_tqf_b),\qquad
 U_D=T^M(\bar f_aY_D),\qquad M\ge1.
}{Q75.0.1}
The exact C066 map, valid for every real Y and every \(a\ge 2,\allowbreak  M\ge 1,\allowbreak \)
is
\fitdisplay{
I_q(T^M(\bar f_aY))^3\le(1-d/72)E|Y|^8,
 \quad I_q(U)=E[q^{-1/3}|U|^{8/3}].
}{Q75.0.2}
Consequently the new \(s\ge 3\) family satisfies
\fitdisplay{
\boxed{I_q(U_D)^3\le(1-d/72)e^{-3d/20}
 \quad(3\le s<a,\ b\ge2,L\ge2,M\ge1).}
}{Q75.T080.4}
The two connector lengths are independent. Section~\ref{sec:completion} gives the exact
arbitrary-joint-law and original P5 applications, with no assumption of
independent branch labels.

This module does not cover unit inner L for \(s\ge 2,\allowbreak \) arbitrary terminal
tables, multiple unequal residual pairs at the root, or every central
role. Other work may supply separate directions in some of those cases;
none is imported here. 

\paragraph{Q75.1 Original objects, positive convolution, and row bounds}\label{part:Q75.1}

Delete zero-mass atoms from the finite probability space and use its expectation and L2 inner product. Set
\fitdisplay{
Jv=(Ev)1,\quad A=J+K,\quad T=J-K=A^*,\quad
 K(x,y)=-K(y,x),\quad |K(x,y)|\le1,
}{}
\fitdisplay{
g=K1,\ p=1+g,\ q=1-g,\ r=Kg,\ s_0=Kr,\ t=Ks_0,
\quad \delta_j=\|K^{j}\one\|_2^2,\quad d=\delta_1,
}{}
\fitdisplay{
f_n=A^{n}\one,\quad \bar f_n=T^{n}\one,\quad
 \theta=13/32,\quad\kappa=19/32,\quad\tau=d-\delta_2.
}{}
The symbol s is an arm length; \(s_0=K^3\) 1 is a function. This avoids
confusing a scalar length with the third score.

The C014 calculus supplies
\fitdisplay{
\|K\|^2\le\theta,\quad0\le d\le1/3,\quad
 d^2\le \delta_2\le\theta d,\quad \delta_{j+1}\le\theta \delta_j,
}{}
\fitdisplay{
\kappa d\le\tau\le d-d^2\le2/9,\quad
 Er=-d,\quad Es_0=0,\quad Et=\delta_2,
\quad Egr=0,\quad Egs_0=-\delta_2,\quad Ert=-\delta_3.
}{Q75.1.1}
All actual path and tree messages are positive on the positive-mass
atoms. In particular p, \(q\ge \mu(x)>0.\) A and T preserve order.

Define the finite coefficients and finite polynomials
\fitdisplay{
\beta_1=\beta_2=0,\qquad
 \beta_j=\sum_{i=1}^{\lfloor(j-1)/2\rfloor}(-1)^{i-1}\delta_i\quad(j\ge3),
\qquad G_n(z)=\sum_{i=0}^{n-1}z^i.
}{}
The original positive finite convolution is
\fitdisplay{
\sum_{j=0}^nK^{j}\one=f_n+\sum_{j=3}^n\beta_j f_{n-j},
 \qquad\tau\le\beta_j\le d\quad(j\ge3).
}{Q75.1.2}
The same finite identity holds after K is replaced everywhere by -K.


For every original vector v, row Cauchy gives \(|Kv(x)|\le \lVert v\rVert _2.\)
The useful two-row estimate is
\fitdisplay{
|K^2v(x)|^2\le\|v\|_2^2/3.
}{Q75.1.3}
This is \eqref{eq:short-one:two-row}; its proof uses conjugated legal skew kernels and the original row measure.

At \(d=0,\allowbreak  g=0\) and every actual tree message here is one by leaf integration,
even when K is nonzero. All theorem boundaries hold directly. Subsequent
divisions by d are restricted to \(d>0.\)

\paragraph{Q75.2 The exact scalar input and its positive majorant interpretation}

Use the positive-mean exponential inequality \eqref{eq:short-one:positive-mgf}. If \(|Z|\le U\) and \(EZ^2\le hU^2,\allowbreak  0\le h<1,\allowbreak \) then, for \(\lambda\ge 0,\allowbreak \)
\fitdisplay{
Ee^{\lambda Z}\le
 \exp\{\lambda\max(EZ,0)+\lambda^2(1+h)^2U^2/8\}.
}{Q75.2.1}
Here all moments belong to the same original finite variable.
A nonnegative upper bound \(M\ge EZ\) may replace max(EZ, 0).

The full proof, allowing either sign of the actual mean, is Lemma~\ref{lem:short-one:exponential}.

At \(\lambda=8/B\) the two forms needed below are:
\fitdisplay{
Z=KV,\ \|V\|^2\le D
 \quad\Longrightarrow\quad
 \log Ee^{8Z/B}\le8M/B+8c_*D/B^2,
 \quad c_*=(1+\theta)^2=2025/1024;
}{Q75.2.2}
\fitdisplay{
Z=K^2V,\ \|V\|^2\le D
 \quad\Longrightarrow\quad
 \log Ee^{8Z/B}\le8M/B+6D/B^2.
}{Q75.2.3}
For (Q75.2.2), the range squared is D and the raw second is at most theta D.
For (Q75.2.3), (Q75.1.3) gives range squared D/3 and \(\theta^2<1/6\) gives raw
second at most D/6, so \(h=1/2\) is legitimate.

Apply AMGM to the following positive factors. With \(W_2=Aq=p-r\) and \(X=T^L(qf_b),\allowbreak \) there is
an actual positive upper bound
\fitdisplay{
0<X\le C[Bq+Br+Z-a],\qquad C,B>0.
}{Q75.2.4}
Thus AMGM between \(W_2\) and the bracket divided by B gives
\fitdisplay{
W_2X\le CB\left(1+\frac{Z-a}{2B}\right)^2.
}{Q75.2.5}
The base is half a sum of nonnegative quantities. Consequently its eighth
power is at most \(C^8B^8\) exp(8(Z-a)/B), including at a zero base.

For a specified nonnegative mean upper M and vector bound D, call the
right-hand logarithmic estimate
\fitdisplay{
\mathcal E=8\log C+8\log B-8a/B+8M/B+
 \begin{cases}6D/B^2,&Z=K^2V,\\8c_*D/B^2,&Z=KV.\end{cases}
}{Q75.2.6}
Section~\ref{proof:Q5-C071} proves every displayed actual mask and
its complete exponent, including the \(b=2,L=2\) row. These complete
expressions are used below for the longer root factors.

\paragraph{Q75.3 The compatible scalar table}\label{part:Q75.3}

\paragraph{Q75.3.1 Exact parameters for every b, L}

For \(b=2\) set \(\gamma=d-\delta_2/2\) and \(B_2=1-\gamma.\) The same original terminal
has \(TF\le q-\gamma-s_0\) and \(T^2F\le B_2q+r+t.\) The C071 finite masks yield:


\begingroup\small
\setlength{\tabcolsep}{3pt}
\setlength{\LTleft}{0pt}\setlength{\LTright}{\fill}
\begin{longtable}{@{}>{\raggedright\arraybackslash}p{0.1200\linewidth}>{\raggedright\arraybackslash}p{0.1600\linewidth}>{\raggedright\arraybackslash}p{0.1600\linewidth}>{\raggedright\arraybackslash}p{0.1600\linewidth}>{\raggedright\arraybackslash}p{0.1600\linewidth}>{\raggedright\arraybackslash}p{0.1600\linewidth}@{}}
\caption{Finite masks and paid scalar bounds for \(b=2\).}\label{tab:parameters:q5-c075:1}\\
\toprule
\(L\) & C & B & a & Z and its row form & Paid scalar upper \\
\midrule\endfirsthead
\toprule
\(L\) & C & B & a & Z and its row form & Paid scalar upper \\
\midrule\endhead
\bottomrule\endfoot
\(L=2\) & 1; & \(B_2\); & 0; & \(K^2(\gamma 1+r)\); & \ensuremath{\le}-d/8 by \mbox{\hyperref[part:Q75.3.2]{part Q75.3.2}}. \\
\addlinespace[2pt]
\(L=3\) & 1; & \(B_2\); & \(\tau\); & \(-K^2(g+s_0)\); & \(\le -3d/4\), by \eqref{eq:ordinary-short2-L3-paid}. \\
\addlinespace[2pt]
\(L=4\) & 1; & \(1-\gamma-\tau\); & \(d(1-\gamma)\); & \(K^2[\tau 1-(1-\gamma)g+r+t]\); & \(\le -d\), by \eqref{eq:ordinary-short2-L4-paid}. \\
\addlinespace[2pt]
\(L=5\) & 1; & \(1-2\tau\); & \(\tau\); & \(K[\tau+zP_3(z)]g,\allowbreak  z=-K\); & \(\le -2d/3\), by \eqref{eq:ordinary-short2-L5-paid}. \\
\addlinespace[2pt]
\(L\ge 6\) & 1; & \(1-2\tau\); & \(\tau\); & \(K[\tau+zP_{L-2}(z)]g\); & \(\le -4d/5\), by \eqref{eq:ordinary-short2-Lge6-paid}. \\
\end{longtable}
\endgroup


Here \(P_n(z)=G_n(z)+ z^{n+1}\). For \(L=2\) one may take \(D=\delta_2\) and
\(M=\delta_2-\gamma d\ge 0.\) For \(L=3\) the exact vector norm is
\(d-2\delta_2+\delta_3\) and \(M=0.\) For \(L=4\) its exact norm is
(\(1-\gamma)^2d+\delta_2-2\delta_3+\delta_4-\tau^2,\allowbreak \) and the used upper mean is
(\(\theta\,\allowbreak \kappa-\tau)d\ge 0.\) The scalar proof uses their stated upper
bounds.
For \(L\ge 5\) the old one-row vector is Qg with \(Q=\delta+zP_n,\allowbreak  \delta=\tau.\) Its old bound \(D=A(\tau)d\) and mean \(M=(\theta-\tau)d\) are
those in Table~\ref{tab:short2:payments}, with
\fitdisplay{
A(\tau)=A_0-2\theta\sigma\tau+\tau^2,
\quad(A_0,\sigma)=
 \begin{cases}(\theta,1),&L=5,\\(\theta-7\theta^2/15,\kappa),&L\ge6.
 \end{cases}
}{Q75.3.1}
The convex endpoint payment there covers the full region
kappa \(d\le \tau\le d-d^2,\allowbreak \) not a finite sample of moments.

For every \(b\ge 3\) use the exact C058/C071 same-direction finite toolkit:
\fitdisplay{
R=\sum_{j=2}^bK^{j}\one=K G_{b-1}(K)g,\quad
 m_b=1-\beta_b/2,\quad k=1/m_b,\quad t_0=\beta_{b+1},
}{}
\fitdisplay{
e=d-\beta_{b+2},\quad V=ER^2,\quad\rho=E(rR),
\quad\gamma=kt_0-k^2V/2,\quad a_1=ke,\quad B_*=\gamma+a_1.
}{Q75.3.2}
In particular \(ER=-t_0,\allowbreak  E(gR)=-e\le 0,\allowbreak  1\le k\le 6/5,\allowbreak  3d/10\le \gamma\le 6d/5,\allowbreak  5d/8\le B_*\le 3d/2,\allowbreak  0\le a_1\le 39d/80,\allowbreak \) and \(m_b^2\le 1-209d/384.\)
Define
\fitdisplay{
P_n(z)=G_n(z)+kz^{n+1}G_{b-1}(-z),\quad z=-K.
}{Q75.3.3}
Every actual upper mask is a finite positive convolution. Specifically,
with \(V_L=\sum _{j=0}^L z^j1+kz^L\) R, its upper bracket before discarding
any nonnegative loss is
\fitdisplay{
X_{b,L}\le m_b^2\left[V_L-\gamma \bar f_{L-1}-a_1\bar f_{L-2}
                  -\sum_{j=3}^L\varepsilon_j\bar f_{L-j}\right],
}{Q75.3.4}
where \(\epsilon_3=d-k \rho\ge 41d/80\) and \(\epsilon_j\ge \tau\) for \(j\ge 4.\)
Only existing indices are used. These exact masks produce (Q75.2.4) with:


\begingroup\small
\setlength{\tabcolsep}{3pt}
\setlength{\LTleft}{0pt}\setlength{\LTright}{\fill}
\begin{longtable}{@{}>{\raggedright\arraybackslash}p{0.1200\linewidth}>{\raggedright\arraybackslash}p{0.1600\linewidth}>{\raggedright\arraybackslash}p{0.1600\linewidth}>{\raggedright\arraybackslash}p{0.1600\linewidth}>{\raggedright\arraybackslash}p{0.1600\linewidth}>{\raggedright\arraybackslash}p{0.1600\linewidth}@{}}
\caption{Finite masks and paid scalar bounds for \(b\ge3\).}\label{tab:parameters:q5-c075:2}\\
\toprule
\(L\) & C & B & a & Z and its row form & Paid scalar upper \\
\midrule\endfirsthead
\toprule
\(L\) & C & B & a & Z and its row form & Paid scalar upper \\
\midrule\endhead
\bottomrule\endfoot
\(L=2\) & \(m_b^2\); & \(1-\gamma\); & \(a_1\); & \(K^2(\gamma 1+kR)\); & \(\le -d/8,\allowbreak  \eqref{eq:short2:longer-L2-paid}\). \\
\addlinespace[2pt]
\(L=3\) & \(m_b^2\); & \(1-B_*\); & \(d-k \rho\); & \(K[\delta+zP_1(z)]g,\allowbreak  \delta=a_1\); & \(\le -d/10,\allowbreak  \eqref{eq:short2:longer-L3-paid}\). \\
\addlinespace[2pt]
\(L=4\) & \(m_b^2\); & \(1-41d/80\); & \(\kappa d\); & \(K[\delta+zP_2(z)]g,\allowbreak  \delta=41d/80\); & \(\le -8d/25,\allowbreak  \eqref{eq:short2:longer-L4-paid}\). \\
\addlinespace[2pt]
\(L\ge 5\) & \(1-d/2\); & \(1-d/2\); & \(d/2\); & \(K[\delta+zP_{L-2}(z)]g,\allowbreak  \delta=d/2\); & \(\le -2d/5,\allowbreak  \eqref{eq:short2:longer-tail-paid}\). \\
\end{longtable}
\endgroup


In the last row \(m_b^2\) is replaced by the larger positive factor 1-d/2
only after the actual upper bracket is known to be nonnegative.
For \(L=2\) the vector squared norm is at most \(k^2V,\allowbreak \) its actual mean is
k rho-gamma d, and \(M=max(k\) rho-gamma d, 0). In the three other
rows \(M=(\theta-\delta)d\ge 0\) and the old squared vector bounds D/d are,
respectively,
\fitdisplay{
351/1280+\delta^2,\qquad\theta(71/63)+\delta^2,
 \qquad\theta+d^2/4.
}{Q75.3.5}
The nonpositive real cross term in those masks is retained in the
original scalar bounds. The finite end sign in the \(L=2\) proof can be
positive for \(b=2\) mod 4; C071 pays it with its same \(\delta_3\) improvement.
We use that complete paid formula, with its actual b, \(m_b,\allowbreak \) R, not a
sign-free relaxation of it.

Every B in the two tables is in (0, 1]. The necessary explicit lower
bounds are \(B_2\ge 13/18,\allowbreak  1-\gamma-\tau\ge 1/2,\allowbreak  1-2\tau\ge 5/9,\allowbreak  1-\gamma\ge 17/25\) in (Q75.3.2), and \(1-B_*\ge 1/2.\)

\paragraph{Q75.3.2 The terminal-two, length-two payment}\label{part:Q75.3.2}
For this actual row \(B=B_2\), \(Z=K^2(\gamma\one+r)\),
\(D=\delta_2\), and \(M=\delta_2-\gamma d\).
Equation~\eqref{eq:short2:L2-paid} proves its complete scalar
upper bound \(\mathcal E\le-d/8\).
In particular every row in the two tables has \(\mathcal E\le-d/10\).

\paragraph{Q75.4 Finite coefficients}\label{part:Q75.4}
For \(z=i\sqrt x\), \(0\le x\le\theta\), the finite coordinates are
\fitdisplay{
G_N(z)=A_N(x)+zB_N(x),\quad1-x\le A_N,B_N\le1\quad(N\ge2).
}{Q75.4.1}
At \(N=1\), \((A_N,B_N)=(1,0)\).
For \(P_n=G_n+z^{n+1}\), or the actual polynomial (Q75.3.3), write
\(P_n=A_P+zB_P\). Equation~\eqref{eq:short-one:coordinates} gives
\fitdisplay{
1-x\le A_P,B_P\le1\quad(n\ge3).
}{Q75.4.2}
Its short rows are
\fitdisplay{
\begin{gathered}
n=1:\quad A_P=1-kxA_N\ge1-6x/5,\quad0\le B_P=kxB_N\le6x/5;\\
n=2:\quad A_P=1-kx^2B_N\ge1-x,\quad
1-6x/5\le B_P=1-kxA_N\le1.
\end{gathered}
}{Q75.4.3}
Thus \(\Re P_n\ge1/2\) at every index used below.
The same finite spectral calculus, Lemma~\ref{lem:odd-common:spectral},
also gives
\fitdisplay{
\|G_m(-K)\|^2\le A_m,\quad A_2=1+\theta,\quad A_m=1\ (m\ge3).
}{Q75.4.4}
The proof in Lemma~\ref{lem:short-one:packet} includes every finite
endpoint and both \(N=1\) and \(N\ge2\).

\paragraph{Q75.5 The actual shortest-three root}\label{part:Q75.5}

The exact root identity is
\fitdisplay{
W_3=A \bar f_2=p-r+s_0-d=W_2+s_0-d.
}{Q75.5.1}
In each positive factorization (Q75.2.4), AMGM with \(W_3\) therefore gives
\fitdisplay{
W_3X\le CB\left(1+\frac{Z+Bs_0-a-Bd}{2B}\right)^2.
}{Q75.5.2}
The base remains nonnegative because the root is actual and the
terminal upper is positive. The new fluctuation has the same mean,
and its deterministic loss adds -8d to (Q75.2.6).

For \(b=2,\allowbreak  L=2,\allowbreak \) the new two-row vector is gamma1+Bg+r. Its squared norm
is \(B^2d+\delta_2-d^2+\delta_2^2/4\le B^2d+\delta_2.\) Thus its cost over the paid expression
of \mbox{\hyperref[part:Q75.3.2]{part Q75.3.2}} increases by at most 6d. The added -8d leaves a bound
at most -\(d/8-2d\le -d.\)
For \(b\ge 3,\allowbreak  L=2,\allowbreak \) the new two-row vector is gamma1+Bg+kR. Since
\(E(gR)=-e\le 0\) and \(Eg=0,\allowbreak \) its squared norm is at most \(k^2V+B^2d.\)
The same argument gives at most -\(d/8-2d\le -d.\) This reuses the explicit
scalar formula with its positive mean, not just the old eighth bound.

For \(b=2,\allowbreak  L=3,\allowbreak \) the new vector is -(gamma \(g+s_0).\) On the original
skew spectrum,
\fitdisplay{
(\gamma-x)^2\le(1-x)^2,
 \quad (1-x)^2-(\gamma-x)^2=(1-\gamma)(1+\gamma-2x)\ge0.
}{}
The two-row norm bound still applies. At \(b=2,\allowbreak  L=4\) the new vector is
tau1-tau g+r+t. Orthogonality of g to 1, r, t and \(\tau<1-\gamma\) show
its squared norm is no larger than the old one. Means are unchanged
in both cases, and their added -8d is more than required.

In every remaining table row, \(Z=K\)[\(\delta+zP_n(z)\)]g with \(z=-K.\)
The new vector is [\(\delta+z(P_n-B)\)]g. For every original frequency,
\fitdisplay{
|\delta+z(P_n-B)|^2-|\delta+zP_n|^2
 =Bx[B-2\Re P_n].
}{Q75.5.3}
This is nonpositive by \mbox{\hyperref[part:Q75.4]{part Q75.4}} and \(0<B\le 1.\) Thus both the original
row range and the theta raw-second estimate use no larger vector
norm. Again the mean is unchanged, and the exponent improves by -8d.
The complete old table therefore proves (Q75.T080.2) for all \(b\ge 2,\allowbreak  L\ge 2.\)
No step asserted \(W_3\le W_2\) pointwise; that assertion would be false.

\paragraph{Q75.6 Every longer shortest arm: an actual loss and a centered perturbation}

Fix the root arm length \(s\ge 4\) and put
\fitdisplay{
m=s-2\ge2,\quad c=1-\tau\ge7/9,\quad
 U=K^2G_m(-K)g,\quad\beta=\beta_s,
\quad H=G_m(-K)g-\tau1.
}{Q75.6.1}
The letter m here is a root-arm polynomial length; it is not \(m_b\)
from the terminal toolkit. Their values are independent.

Apply positive A to the negative-sign version of (Q75.1.2) for \(\bar f_{s-1}.\)
The left side becomes \(p-r+U-\beta_s,\allowbreak \) while the last convolution
term on the right is \(\beta_{s-1}p.\) Every other term is nonnegative.
Since \(s-1\ge 3,\allowbreak \) retaining \(\beta_{s-1}p\ge \tau\) p gives
\fitdisplay{
0<W_s\le cp-r+U-\beta.
}{Q75.6.2}
The upper side is nonnegative because it is the actual \(W_s\) plus
the remaining nonnegative convolution terms. In particular, it is
not an arbitrarily signed formal series.

Combine (Q75.6.2) with the positive terminal upper (Q75.2.4). AMGM between
the root upper divided by c and the terminal upper divided by CB
gives, with \(h=B/c,\allowbreak \)
\fitdisplay{
W_sX\le CBc\left(1+
 \frac{Z+hK^2H-a-h\beta}{2B}\right)^2.
}{Q75.6.3}
To check the algebra, the new fluctuation is
Z+(B/c)(U-tau r), and \(K^2H=U-\tau\) r. The deterministic addition
is -(B/c)beta; the constant prefactor is CBc. Both AMGM factors
are nonnegative. This scaling is needed: it retains the actual
loss tau p without introducing a fictitious root factor.

\paragraph{Q75.6.1 The same original perturbation has a nonpositive mean}

Let beta'=E \(G_m(-K)g.\) The finite skew identities give
\fitdisplay{
\beta'=\beta_{m+1},\quad
 \tau\le\beta'\le d,\quad 0\le EH=\beta'-\tau\le \delta_2,
}{}
\fitdisplay{
\|H\|^2=\|G_m(-K)g\|^2-2\tau\beta'+\tau^2
 \le A_m d-\tau^2\le A_m d.
}{Q75.6.4}
To control its \(K^2\) mean without separating incompatible moment extrema,
apply the original norm bound to the centered vector r+d1:
\fitdisplay{
\delta_3-2d\delta_2+d^3=\|K(r+\delta_1)\|^2\le\theta(\delta_2-d^2).
}{}
It follows that
\fitdisplay{
\delta_2-\delta_3-\tau d\ge(\kappa-d)(\delta_2-d^2)\ge0.
}{Q75.6.5}
For every \(m\ge 2,\allowbreak \) the finite alternating mean of U begins -\(\delta_2\) and
is at most -\(\delta_2+\delta_3.\) Hence the same H satisfies
\fitdisplay{
EK^2H=EU+\tau d\le0.
}{Q75.6.6}
Thus adding \(hK^2H\) never increases any nonnegative mean upper M used
in the old scalar table. This is also valid when an actual mean was
already negative.

\paragraph{Q75.6.2 A nonpositive covariance for every terminal at \(L=2\)}

The old two-row vector for \(L=2\) is \(V_0=gamma1+kR.\) For \(b\ge 3\) its
parameters are (Q75.3.2). For \(b=2\) use \(R=r,\allowbreak  k=1,\allowbreak  t_0=d,\allowbreak  V=\delta_2\) and the
same identity \(\gamma=kt_0-k^2V/2.\) In both cases
\fitdisplay{
EV_0=-k^2V/2\le0.
}{}
We prove
\fitdisplay{
\operatorname{Cov}(R,G_m(-K)g)\le0,
 \qquad E(V_0H)\le0.
}{Q75.6.7}
The second follows from the first and \(EH\ge 0.\)

For \(b\ge 3\) put \(N=b-1\ge 2,\allowbreak \) and at an original skew rotation block write
\fitdisplay{
G_N(z)=A+zB_c,\qquad G_m(z)=C+zD,\qquad z=i\sqrt x.
}{}
All four real coefficients belong to [1-x, 1], hence to [kappa, 1].
Let \(\nu_1\) be the spectral measure of the original vector 1 for -\(K^2.\)
It has mass one, including any zero block. The corresponding measure
of \(g=K1\) is x \(\nu_1.\) This is spectral bookkeeping in the original
finite Hilbert space; it is not a changed host or a new row bound.
Skew adjunction gives the exact three quantities
\fitdisplay{
t_0=\int xA\,d\nu_1,\qquad
 \beta'=\int xD\,d\nu_1,
\quad E[R G_m(-K)g]=-\int x^2(AD+B_cC)\,d\nu_1.
}{Q75.6.8}
Put \(J_0=\int x^2AD\,d\nu_1\). Pointwise \(B_cC\ge \kappa^2AD,\allowbreak \) and
A/D lies in [kappa, 1/kappa]. By \(2uv\le u^2+v^2\) and Cauchy with
the probability measure \(\nu_1,\allowbreak \)
\fitdisplay{
t_0\beta'\le\tfrac12(t_0^2+\beta'^2)
 \le\tfrac12\int x^2(A^2+D^2)\,d\nu_1
 \le\frac{\kappa+\kappa^{-1}}2 J_0
 \le(1+\kappa^2)J_0.
}{Q75.6.9}
The final comparison is precisely \(2\kappa\ge 1.\) Combining (Q75.6.8)-(Q75.6.9)
proves the first inequality of (Q75.6.7).


Here the covariance in (Q75.6.7) equals
\fitdisplay{
-\operatorname{Cov}_{\nu_1}(x,xD(x)).
}{}
The function xD(x) is a finite alternating sum \(x-x^2+x^3-...\) .
It is increasing on [0, theta]. Indeed its derivative begins \(1-2x>0\);
each subsequent positive/negative pair
(\(2j+1)x^{2j}-(2j+2)x^{2j+1},\allowbreak  j\ge 1,\allowbreak \) is nonnegative, and an
unpaired last positive term is nonnegative. For the constant \(D=1\)
the assertion is immediate. The finite double-integral identity
\fitdisplay{
2\operatorname{Cov}_{\nu_1}(u,v)
 =\iint[u(x)-u(y)][v(x)-v(y)]\,d\nu_1(x)d\nu_1(y)
}{}
therefore proves the needed sign. It includes the zero block.

Finally, for both terminal ranges,
\fitdisplay{
E(V_0H)=k\operatorname{Cov}(R,G_m(-K)g)+(EV_0)(EH)\le0.
}{}
Thus the new two-row vector \(V_0+hH\) obeys
\fitdisplay{
\|V_0+hH\|^2\le\|V_0\|^2+h^2A_m d.
}{Q75.6.10}
The old relaxed norm D from \mbox{\hyperref[part:Q75.3]{part Q75.3}} may replace \(\lVert V_0\rVert ^2\) on
the right. The underlying vector and its cross term remain actual.

\paragraph{Q75.6.3 The other two-row cases \(b=2,\allowbreak  L=3,\allowbreak \) 4}

For \(L=3,\allowbreak  V_0=-(g+s_0),\allowbreak  EV_0=0.\) Its product with \(G_m(-K)g\) has
spectral real part -(1-x)C, so \(E(V_0H)\le 0.\)

For \(L=4,\allowbreak  V_0=tau1-(1-\gamma)g+r+t\) and \(EV_0=0.\) The product with
H, using the original spectral measure \(\nu_g,\allowbreak \) is exactly
\fitdisplay{
E(V_0H)=\tau\beta'-(1-\gamma)\int C\,d\nu_g
                  -\int x(1-x)D\,d\nu_g
 \le[\tau-\kappa(1-\gamma)]d\le0.
}{Q75.6.11}
For the last sign use \(\tau\le 2/9\) and
\(\kappa(1-\gamma)\ge \kappa(13/18)=247/576>2/9.\)
The same bound (Q75.6.10) therefore holds in every old two-row case.

\paragraph{Q75.6.4 Every remaining one-row case}

The old vector is \(V_0=Q(-K)g\) with \(Q=\delta+zP_n(z).\)
The added vector in the one-row representation is hKH, since
\(Z+hK^2H=K(V_0+hKH).\) At \(z=-K\) we have
\fitdisplay{
KH=-(\tau+zG_m(z))g.
}{}
We prove its cross with \(V_0\) is nonpositive at every original frequency.
Write \(G_m=C+zD\) and \(P_n=A_P+zB_P\) as in \mbox{\hyperref[part:Q75.4]{part Q75.4}}. Direct multiplication
gives
\fitdisplay{
\Re\{Q\overline{\tau+zG_m}\}
 =\delta\tau+x[\Re(P_n\overline{G_m})-\delta D-\tau B_P].
}{Q75.6.12}
Here \(\delta\ge 0.\) Every table row with \(n\ge 3\) has \(\delta\le 2/9.\)
By (Q75.4.1)-(Q75.4.2) the square bracket is at least
\fitdisplay{
(1+x)(1-x)^2-\delta-\tau
 \ge16245/32768-4/9>0.
}{Q75.6.13}
The function (\(1+x)(1-x)^2\) decreases on [0, theta]; its endpoint is
(\(45/32)(19/32)^2.\) The last rational comparison is
\(146205>131072.\) This covers every \(b=2,\allowbreak  L\ge 5\) and \(b\ge 3,\allowbreak  L\ge 5.\)

For \(b\ge 3,\allowbreak  L=4,\allowbreak  n=2\) and \(\delta=41d/80\le 2/9.\) Use (Q75.4.3) to bound
the square bracket by
\fitdisplay{
(1-x)(1-6x^2/5)-\delta-\tau
 \ge78014/163840-4/9>0.
}{Q75.6.14}
The displayed frequency function decreases throughout the interval:
its derivative is \(-1-12x/5+18x^2/5<0\) for \(0\le x\le\theta\). At theta it
has the stated value; the final comparison is \(702126>655360.\)

For \(b\ge 3,\allowbreak  L=3,\allowbreak  n=1\) and \(\delta=a_1\le 13/80.\) Again (Q75.4.3) gives
\fitdisplay{
\Re(P_n\overline{G_m})-\delta D-\tau B_P
 \ge(1-6x/5)(1-x)-\delta-(6/5)\tau x
 \ge779/2560-13/80-13/120=257/7680>0.
}{Q75.6.15}
Both negative replacements are legitimate since \(D\le 1\) and
\(B_P\le 6x/5.\) The first product decreases on [0, theta].

These are the full possible n-ranges, including \(n=1,\allowbreak \) 2. Equations
(Q75.6.12)-(Q75.6.15) show \(E(V_0 KH)\le 0.\) Therefore
\fitdisplay{
\|V_0+hKH\|^2\le\|V_0\|^2+h^2\|KH\|^2
 \le\|V_0\|^2+h^2\theta A_m d.
}{Q75.6.16}
This norm bound uses
\(\lVert H\rVert ^2\) from (Q75.6.4). Row Cauchy and the spectral theta bound apply to this vector.

\paragraph{Q75.7 Payment of every longer root, on the whole continuous domain}

All old mean uppers stay valid by (Q75.6.6). In a two-row case, (Q75.6.10)
adds at most \(6h^2A_m d/B^2=6A_m d/c^2\) to its exponential cost.
In a one-row case, (Q75.6.16) adds at most \(8c_*\theta A_m d/c^2.\)
Since \(6<8c_*\theta=26325/4096,\allowbreak \) both are paid by the common upper
increment
\fitdisplay{
\Delta\le8\log c-8\beta/c+8c_*\theta A_m d/c^2.
}{Q75.7.1}
The first two terms come from the actual factor c and actual loss
beta in (Q75.6.3). 

For \(m\ge 3,\allowbreak \) equivalently \(s\ge 5,\allowbreak  A_m=1\) and \(\beta\ge \tau.\) Using
log \(c\le -\tau,\allowbreak  \tau/d\ge \kappa\) and \(c\ge 7/9,\allowbreak \) divide by \(8d/c^2\) to get
\fitdisplay{
\frac{c^2\Delta}{8d}
 \le-\kappa(c^2+c)+c_*\theta
 \le-\frac{133}{162}+\frac{26325}{32768}
 <-\frac1{64}.
}{Q75.7.2}
For the last comparison, \(133/162-26837/32768\) equals
\(10550/5308416>0.\) Thus \(\Delta\le -d/8\) in this range.

For \(m=2,\allowbreak \) equivalently \(s=4,\allowbreak \) the actual \(\beta_4\) is d and
\(A_m=1+\theta.\) Consequently
\fitdisplay{
\frac{c^2\Delta}{8d}
 \le-\kappa c^2-c+c_*\theta(1+\theta)
 \le-\frac{2947}{2592}+\frac{1184625}{1048576}
 <-\frac1{160}.
}{Q75.7.3}
Indeed \(1184625/1048576<113/100\) and
\(2947/2592-113/100=451/64800>1/160.\) These are exact rational comparisons;
no continuous-domain optimization is used.
Thus in every \(s\ge 4\) case \(\Delta\le -d/20.\)

The complete old table, including \mbox{\hyperref[part:Q75.3.2]{part Q75.3.2}}, has scalar exponent
at most -d/10. Adding the fully paid increment proves
\fitdisplay{
E(W_sX_{b,L})^8\le e^{-3d/20}
 \qquad(s\ge4,b\ge2,L\ge2).
}{}
\mbox{\hyperref[part:Q75.5]{part Q75.5}} supplies \(s=3.\) Finally
\fitdisplay{
e^{-3d/20}\le\frac1{1+3d/20}
 =1-\frac{3d/20}{1+3d/20}\le1-d/7
 \qquad(0\le d\le1/3).
}{}
This proves (Q75.T080.1) in every integer and continuous parameter,
with no exceptional intermediate d interval.
